\section{ZIPs as Research Proposals}

\label{sec:zips}
%----------------------------------------------------------------------

\subsection{From Improvement Proposals to Research Proposals}

Zoo Improvement Proposals (\ZIP{}s) \cite{zoo-dao-governance} are the governance mechanism for the \Zoo{} Network. We extend \ZIP{}s to serve as research funding proposals, replacing traditional grant applications:

\begin{definition}[Research ZIP]
A Research \ZIP{} $Z = (\textit{title}, \textit{abstract}, \textit{team}, \textit{budget}, \textit{milestones}, \textit{timeline}, \textit{capsule\_id})$ is a proposal for community-funded research, submitted to \GChain{} for \DAO{} vote.
\end{definition}

\subsection{Proposal Lifecycle}

\begin{enumerate}
    \item \textbf{Submission.} Researcher submits a Research \ZIP{} to \GChain{} with a deposit of 10M $\KEEPER$ tokens (refundable upon completion).
    \item \textbf{Community Review.} A 14-day review period where any $\KEEPER$ holder can comment, ask questions, and signal support or opposition.
    \item \textbf{Peer Review.} Three reviewers are selected from the reviewer pool (staked researchers with verified credentials on \IChain{}). Reviewers stake 1M $\KEEPER$ on their review quality.
    \item \textbf{Vote.} A 7-day voting period on \GChain{}. Standard proposals require 51\% majority with 5\% quorum; large grants ($>$100M $\KEEPER$) require 66\% supermajority with 10\% quorum.
    \item \textbf{Funding.} Approved proposals receive milestone-based funding from the \DAO{} treasury.
    \item \textbf{Milestones.} Each milestone requires a deliverable (data, code, paper) committed to the Research Capsule. Milestone completion is verified by the original reviewers.
    \item \textbf{Completion.} Final milestone triggers deposit refund and completion badge on the researcher's \IChain{} \DID{}.
\end{enumerate}

\subsection{Incentivized Peer Review}

Peer review on \Zoo{} is not volunteer labor---it is staked, compensated, and reputation-building:

\begin{itemize}
    \item \textbf{Reviewer Stake.} Reviewers stake 1M $\KEEPER$ when accepting a review assignment. The stake is slashed if the review is flagged as low-quality by the community (>10\% of voters object).
    \item \textbf{Compensation.} Reviewers earn 500K $\KEEPER$ per completed review, paid from the proposal's budget.
    \item \textbf{Reputation.} Completed reviews accrue reputation tokens on the reviewer's \DID{}. High-reputation reviewers are prioritized for future assignments.
    \item \textbf{Conflict Checks.} \IChain{} \DID{} credentials are used to verify that reviewers have no institutional or co-authorship conflicts with the proposer.
\end{itemize}

\subsection{Advantages Over Traditional Grants}

\begin{table}[h]
\centering
\begin{tabular}{lcc}
\toprule
\textbf{Property} & \textbf{Traditional Grants} & \textbf{Research ZIPs} \\
\midrule
Decision timeline & 12--18 months & 21 days \\
Reviewer compensation & \$0 & 500K $\KEEPER$ \\
Review transparency & Opaque & Public \\
Funding source & Government/foundation & Community \DAO{} \\
Geographic restriction & Yes (citizenship, residency) & No \\
Progress tracking & Annual reports & On-chain milestones \\
Reproducibility requirement & Rare & Mandatory \\
\bottomrule
\end{tabular}
\caption{Research ZIPs vs.\ traditional grant funding}
\label{tab:grants_comparison}
\end{table}

%----------------------------------------------------------------------
